% En esta seccion se expresan sus conclusiones del trabajo realizado de manera concreta y coherente.

\section{Conclusiones}

\subsection{Basilio Illescas Andrés}
Los objetivos de esta entrega se cumplieron: la pierna articulada quedó construida y deformando correctamente, el personaje de Mixamo quedó animado, importado y con su jerarquía caracterizada hueso por hueso, y las relaciones padre-hijo quedaron representadas en un diagrama que parte del nodo raíz. Lo que realmente me llevo de la práctica es haber podido comparar, dentro del mismo reporte, una jerarquía construida por nosotros contra una construida por una herramienta profesional. La primera tiene cinco huesos en una sola cadena y la segunda tiene sesenta y siete repartidos en doce niveles, pero al medirlas resultó que la diferencia entre ambas es únicamente de escala: en las dos, cada hueso guarda su transformación respecto de su padre y la posición en el mundo sale de multiplicar la cadena completa desde la raíz. Ese era el concepto de la práctica y verlo sostenerse en los dos extremos me lo dejó mucho más claro que cualquier definición.

La lección de método vino de no conformarme con la afirmación teórica. Que Blender aplique cinemática directa es algo que uno repite con facilidad, pero reconstruir la matriz global de los sesenta y siete huesos a partir exclusivamente de la de su padre y comprobar que la discrepancia máxima es de $4.578 \times 10^{-5}$ convierte la afirmación en una medición. Y el resultado tuvo un detalle que no esperaba: el peor error no cayó en cualquier lugar, sino en una falange, es decir, en uno de los nodos más profundos del árbol. El error numérico se acumula siguiendo exactamente la misma cadena por la que se propaga la transformación, lo cual es evidente en retrospectiva pero no se me había ocurrido antes de verlo.

Si en la Práctica 03 me quedé con que «se ve bien» no equivale a «está bien», aquí me queda una versión complementaria: que dos cosas se vean iguales tampoco significa que lo sean. Los dos archivos de la variación experimental pesan exactamente los mismos \texttt{28\,426\,960} bytes y contienen la misma malla, los mismos huesos y el mismo número de fotogramas; lo único que los distingue son los valores de una curva, la del único hueso que no tiene padre. Que anular esa sola curva sea suficiente para que el personaje deje de avanzar por el mundo, y que al mismo tiempo su balanceo lateral y vertical se conserven idénticos hasta la última cifra, es la demostración más limpia que vi en toda la práctica de qué significa realmente ser el nodo raíz de una jerarquía.

\subsection{Cruz Macedo Samuel Santiago}
Entrando a esta práctica pensé que la parte difícil iba a ser la Actividad 1, porque ahí hay que modelar y colocar los huesos uno mismo, y que recibir un personaje ya rigueado de Mixamo sería el trámite fácil. Terminó pasando lo de siempre: la dificultad no desapareció, nada más se movió de lugar. Construir una cadena de cinco huesos es laborioso pero uno sabe en todo momento qué hay adentro; leer una estructura ajena de sesenta y siete huesos es lo contrario, porque el archivo ya viene resuelto y el trabajo consiste en averiguar con qué criterio lo resolvieron. Es la diferencia entre escribir un programa y tener que entender el de alguien más.

Lo que más me sorprendió fue descubrir que en el esqueleto de Mixamo la mayoría de los huesos están emparentados pero \textbf{no} conectados: cuarenta y siete de los sesenta y siete. En la Actividad 1 habíamos dejado la pierna con todos sus huesos encadenados justamente para que no se separaran al rotar, así que al principio pensé que el archivo de Mixamo traía un defecto. Al revisarlo con calma se entiende que no es un descuido sino una decisión: un hombro no nace en la punta de la columna ni un muslo en la punta de la cadera, y conectarlos obligaría a que la cabeza del hijo coincidiera con la cola del padre. Eso me hizo separar dos ideas que yo tenía mezcladas, porque emparentar y conectar no son lo mismo: lo que hereda las transformaciones es el emparentado, y conectar es apenas una restricción de posición encima.

De la Actividad 3 me quedo con algo que no venía en el manual y que tampoco esperaba que fuera un problema. Al generar el diagrama completo con los sesenta y siete huesos salió una figura de proporción ocho a uno, tan alargada que los nombres quedaban impresos en un tamaño imposible de leer. Tuvimos que probar varias salidas hasta darnos cuenta de que cuarenta de esos huesos son falanges repetidas en diez cadenas idénticas, y que condensarlas en un nodo resumen no pierde información, solo deja de repetirla. Representar una jerarquía, entonces, no es únicamente volcar todos los nodos en una hoja: hay que decidir qué se muestra y qué se resume, igual que cuando en la práctica anterior había que decidir con cuántos lados aproximar un círculo. Considero que los objetivos se cumplieron, y me queda claro que esto es la base de lo que sigue, porque una vez que la cadena de transformaciones está bien armada, la cinemática inversa no es más que el problema de recorrerla en sentido contrario.

\subsection{Medel Piña Ángel de Jesús}
% Conclusión individual del integrante...

\subsection{Pérez Olvera Alexis Abraham}
Esta práctica representó un cambio de paradigma respecto a lo que veníamos trabajando en las entregas anteriores. En las primeras prácticas los objetos eran cuerpos rígidos que se transformaban en bloque mediante una única matriz de modelo; aquí, en cambio, el reto fue entender cómo una geometría continua y orgánica puede deformarse de manera creíble a partir de una estructura interna. Ver cómo la rotación se transmite en cascada desde la cadera hasta las falanges del pie me ayudó a aterrizar por fin la teoría de las estructuras jerárquicas: no hace falta recalcular la posición en el mundo de cada segmento, porque la propia multiplicación acumulada de matrices ($M_{\text{padre}} \cdot M_{\text{local}}$) se encarga de arrastrar a los nodos hijos de forma natural.

El aspecto que más me llamó la atención fue descubrir que el éxito de una armadura no depende únicamente de cómo se coloquen los huesos, sino de la topología previa de la malla. Al inicio pensé que el algoritmo de pesos automáticos resolvería cualquier geometría, pero al probar las primeras flexiones quedó claro que sin bucles de soporte limpios en las zonas de articulación (como la rótula y el pliegue poplíteo), la malla simplemente colapsa y pierde volumen. Igualmente ilustrativo fue el detalle del eje de flexión en la rodilla: colocar los huesos en una línea vertical perfectamente recta parecía lo más intuitivo, pero matemáticamente creaba una ambigüedad que desorientaba la rotación; desplazar la articulación apenas unos milímetros hacia adelante en el eje $Y$ para darle una pre-flexión anatómica solucionó el problema de inmediato. En conclusión, me quedo con la lección de que en animación y modelado jerárquico la geometría, la cinemática y la anatomía tienen que estar pensadas en conjunto desde el primer trazo, ya que ningún rig puede compensar una mala distribución de polígonos.

\subsection{Segura Cedeño Luisa María}
% Conclusión individual del integrante...
